% Grado 11 -- generado por tools/build_tex.py a partir de raw/grado_11.txt
\section{Grado 11}\label{grado:11}

% PDF p. 37
\dba{1}{Comprende la naturaleza de la propagación del sonido y de la luz como fenómenos ondulatorios (ondas mecánicas y electromagnéticas, respectivamente).}

\evidencias

\begin{itemize}
  \item Clasifica las ondas de luz y sonido según el medio de propagación (mecánicas y electromagnéticas) y la dirección de la oscilación (longitudinales y transversales).
  \item Aplica las leyes y principios del movimiento ondulatorio (ley de reflexión, de refracción y principio de Huygens) para predecir el comportamiento de una onda y los hace visibles en casos prácticos, al incluir cambio de medio de propagación.
  \item Explica los fenómenos ondulatorios de sonido y luz en casos prácticos (reflexión, refracción, interferencia, difracción, polarización).
  \item Explica las cualidades del sonido (tono, intensidad, audibilidad) y de la luz (color y visibilidad) a partir de las características del fenómeno ondulatorio (longitud de onda, frecuencia, amplitud).
\end{itemize}

\ejemplo

Dadas las situaciones que se presentan en las fotos, explica, a partir de principios y leyes, lo que observa.

\fig{g11_p37_1}{Ejemplo: (1) un puntero láser apunta hacia un vaso con agua y el haz se desvía al entrar en el agua (refracción de la luz); (2) un lápiz dentro de un vaso con agua se ve quebrado en la superficie (refracción); (3) «Espejos de sonido»: dos reflectores parabólicos enfrentados, un emisor y un receptor, con los frentes de onda del sonido que se reflejan de uno a otro (reflexión del sonido).}


% PDF p. 37
\dba{2}{Comprende que la interacción de las cargas en reposo genera fuerzas eléctricas y que cuando las cargas están en movimiento genera fuerzas magnéticas.}

\evidencias

\begin{itemize}
  \item Identifica el tipo de carga eléctrica (positiva o negativa) que adquiere un material cuando se somete a procedimientos de fricción o contacto.
  \item Reconoce que las fuerzas eléctricas y magnéticas pueden ser de atracción y repulsión, mientras que las gravitacionales solo generan efectos de atracción.
  \item Construye y explica el funcionamiento de un electroimán.
\end{itemize}

\ejemplo

\fig{g11_p37_2}{Ejemplo (solo imagen): circuito sobre una tabla con una pila, un bombillo encendido y una brújula junto al cable. Al circular la corriente, la aguja de la brújula se desvía (experimento de Oersted): una carga en movimiento genera un campo magnético.}


% PDF p. 38
\dba{3}{Comprende las relaciones entre corriente y voltaje en circuitos resistivos sencillos en serie, en paralelo y mixtos.}

\evidencias

\begin{itemize}
  \item Determina las corrientes y los voltajes en elementos resistivos de un circuito eléctrico utilizando la ley de Ohm.
  \item Identifica configuraciones en serie, en paralelo y mixtas en diferentes circuitos representados en esquemas.
  \item Identifica características de circuitos en serie y paralelo a partir de la construcción de circuitos con resistencias.
  \item Predice los cambios de iluminación en bombillos resistivos en un circuito al alterarlo (eliminar o agregar componentes en diferentes lugares).
\end{itemize}

\ejemplo

Dado el siguiente circuito:

\fig{g11_p38_1}{Circuito mixto alimentado por una pila de 4,5 V con cinco bombillos resistivos. Rama superior: dos bombillos en paralelo (uno de ellos marcado X) en serie con un tercero. Rama inferior: dos bombillos en serie. Las dos ramas están en paralelo con la pila. Se pregunta cómo cambia la iluminación de los demás bombillos si se elimina X.}

Identifica los cambios de iluminación en los bombillos restantes si el bombillo (X) se elimina.


